\newpage
\begin{table}[!t]
\centering
\caption{Summary of Core Empirical Findings and Confidence Levels}
\label{tab:summary_findings}
\scriptsize
\renewcommand{\arraystretch}{1.02}
\begin{tabular*}{\columnwidth}{@{\extracolsep{\fill}}p{2.6cm}p{3.8cm}c}
\toprule
\textbf{Empirical Finding} & \textbf{Quantitative Evidence} & \textbf{Confidence} \\
\midrule
\textbf{SysV MQ Low-Payload Lead} & $123.33\,\mu\text{s}$ (64 B), $115.13\,\mu\text{s}$ (256 B); lowest sample mean ($p \ge 0.05$, paired $t$-test; descriptive lead under KVM floor). & Descriptive \\
\addlinespace[1.5pt]
\textbf{Unix Socket Scaling at 256 KB} & $449.64\,\mu\text{s}$ vs. $1144.22\,\mu\text{s}$ for Pipe ($2.54\times$ speedup, paired $t$-test $p < 0.001$, 1,304.4 MB/s wire throughput). & High \\
\addlinespace[1.5pt]
\textbf{Pipe Buffer Saturation Penalty} & Voluntary context switches surge from 20.8k to 146.4k at 256 KB due to 64 KB capacity. & High \\
\addlinespace[1.5pt]
\textbf{Shared Memory Kernel Bypass} & 0.00 data-path syscalls; 2.06 futex calls/exch vs. 4.00 for streams (48.5\% kernel entry drop). & High \\
\addlinespace[1.5pt]
\textbf{KVM Sched. Floor Bounds IPC} & Latencies tightly cluster at 115--140 $\mu$s across transports despite 0.72 $\mu$s driver floor. & High \\
\addlinespace[1.5pt]
\textbf{Static Selector Efficiency} & 0.66 ns evaluation; 71.43\% candidate accuracy (SHM vs. Pipe); 27.68\% mean regret. & High \\
\addlinespace[1.5pt]
\textbf{Offline 3-Tier Policy Capture} & 100.0\% oracle match (7 of 7 payload classes) across evaluated single-pair workloads. & Supported \\
\bottomrule
\end{tabular*}
\vskip 1.5pt
{\tiny\noindent\emph{Confidence taxonomy}: \textbf{High}: Statistically significant separation ($\alpha = 0.05$) or deterministic OS telemetry. \textbf{Descriptive}: Consistent sample-mean ranking without statistically significant separation ($p \ge 0.05$). \textbf{Supported}: Empirically calibrated structural model.\par}
\end{table}

\balance
\section{Conclusion and Future Work}
\label{sec:conclusion}

This research presented \textbf{Mini-A-IPC}, an empirical evaluation platform and adaptive communication framework for Linux IPC. Across over 3.4 million scheduled round-trip exchanges (2.73M valid timed exchanges, $>4.5$M total evaluated transmissions) with environmental isolation, randomized blocking, and CPU pinning on Fedora 45 Linux (kernel 7.3.4, KVM), three core findings were established:
(1)~\textbf{Payload Crossovers:} System V Message Queues deliver lowest small-payload sample latency ($123.33\,\mu\text{s}$ at 64~B, $115.13\,\mu\text{s}$ at 256~B, descriptive lead under the KVM floor), Named FIFOs and Anonymous Pipes dominate intermediate payloads ($119.44\,\mu\text{s}$ at 4~KB), and Unix Domain Sockets outperform pipes by 2.54$\times$ for 256~KB transfers ($449.64\,\mu\text{s}$, 1,304~MB/s, paired $t$-test $p < 0.001$) by preventing buffer exhaustion.
(2)~\textbf{Adaptive Selection:} A static size-threshold selector evaluates in 0.66~ns, achieving 71.43\% candidate accuracy between shared memory and pipes with 27.68\% mean regret, while offline 3-tier calibration delineates optimal static routing and an EWMA telemetry controller provides the architecture for dynamic workloads.
(3)~\textbf{Kernel vs. Scheduler Dynamics:} A custom character device mailbox (\texttt{/dev/aipc}) achieves an uncontended 0.72~$\mu$s loopback floor, demonstrating that hypervisor vCPU scheduling (95--130~$\mu$s) bounds multi-process exchange speed.

\textbf{Future Work:} Directions include bare-metal hardware benchmarking, NUMA-aware allocation, dynamic lock-free ring resizing, an asynchronous \texttt{io\_uring} driver, and live online EWMA evaluation under dynamic burst traffic.
